\section{Background and Related Work}
\label{sec:background}

\subsection{Knowledge-based temporal abstraction}

Knowledge-based temporal abstraction (KBTA) \citep{shahar1997framework}
converts raw point observations into interval-level states, gradients
(trends), events, and contextual periods by applying declarative,
potentially context-sensitive medical domain knowledge. Every KBTA token
is a named clinical concept---for example \emph{decreasing-glucose trend}
over three days, \emph{moderate-anemia}, or \emph{on-basal-insulin
context}---drawn from a curated ontology rather than an unnamed
discretized bin index. This naming makes downstream attribution resolve
to clinical concepts by construction. Interval sequences produced this way
have been mined for outcome-associated patterns using KarmaLego
\citep{moskovitch2015karmalego} with pattern-frequency summarization via
Temporal Probabilistic Profiles \citep{sheetrit2019tpf}, or embedded into
sequence models that consume the interval stream directly. In this paper
we use the Mediator engine \citep{mediator-github}, which applies a curated
knowledge base to emit interval-typed abstractions and, in the same pass,
computes fuzzy-logic compliance scores for each formalized guideline
recommendation \citep{oded2027kineret-jbi}.

\subsection{Deep learning for irregular clinical streams}

Standard approaches to irregular EHR data include GRU-D
\citep{che2018recurrent}, which augments a gated recurrent network with
learned exponential decays that model elapsed time between observations
and the effects of missingness, and STraTS \citep{tipirneni2022strats},
which represents each observation as a (concept, time, value) triplet and
processes the resulting set with a Transformer. Visit-level or
concept-level Transformer variants such as BEHRT \citep{li2020behrt},
TransformEHR \citep{yang2023transformehr} and HEART
\citep{huang2024heart} target longitudinal histories via
diagnosis-code sequences with time-gap embeddings. In all these
architectures the model reads either raw or highly aggregated inputs,
and attributions resolve to raw variables or diagnosis codes rather than
to named clinical concepts.

\subsection{Interval-based Transformers over KBTA}

INTERVenE \citep{oded2026intervene} tokenizes KBTA intervals as pairs of
START/END points, each expanded through four hierarchical embedding
layers---raw concept, TAK abstraction concept, concept-with-value, and
positional role---so that concepts sharing a common ancestor
(\emph{glucose-state} vs \emph{glucose-trend}) or a common interval share
part of their embedding by construction. Continuous time is injected
through a Time2Vec \citep{kazemi2019time2vec} channel that is concatenated
with the four concept channels and linearly projected to the model
dimension, and static patient context is injected through AdaLN-Zero
\citep{peebles2023adaln0}. Two variants share this backbone: an
autoregressive decoder (INTERVenE-Ar) that generates full abstraction
trajectories with a per-step risk readout, and a bidirectional encoder
(INTERVenE-Enc) that yields single-pass joint risk and time-to-event
predictions. On $57{,}078$ MIMIC-IV diabetic admissions, INTERVenE-Enc
reached support-weighted AUPRC $0.672$ against $0.631$ for the strongest
neural baseline (ss-STraTS) with non-overlapping bootstrap intervals; a
concurrent ablation showed that replacing KBTA intervals with a
knowledge-free $\sigma$-bin discretization left the AUPRC essentially
unchanged, so the accuracy gain on that cohort was attributed to the
Transformer over persistent intervals rather than to the semantic content
of the intervals themselves.

The present paper uses INTERVenE-Enc only. INTERVenE-Ar is a generative
inspection tool rather than a discriminative predictor and is scored
under a different contract; it is not part of the ladder here. Sparse-
transcoder attribution \citep{dunefsky2024transcoders,rajamanoharan2024jumping},
originally developed for INTERVenE-Enc on MIMIC-IV, is used unchanged in
this study to interpret which tokens the model relies on when the QA
signal is supplied.

\subsection{Guideline adherence as a temporal-abstraction problem}

Assessing adherence to a clinical guideline at the level of its individual
recommendations, over the course of an admission, is naturally a temporal
problem: many recommendations concern a decision that must be made
\emph{whenever} some clinical state obtains (``insulin following
hyperglycemia''), rather than a discrete action taken at admission
(``obtain HbA1c on admission''). Rule- and pattern-based frameworks over
KBTA intervals \citep{oded2027kineret-jbi} formalize each recommendation
as a temporal pattern with a fuzzy-logic compliance score, yielding a
continuous score wherever the pattern applies and a null result where it
does not. The Kineret adherence study measured such scores for fourteen
recommendations of the Israeli national inpatient hyperglycemia guideline
\citep{hochberg2023inpatient} across $48{,}354$ admissions and identified,
after propensity weighting, above-median adherence as associated with
lower odds of $30$-day mortality (OR $=0.557$) and severe hyperglycemia
(OR $=0.485$); the recommendations most strongly linked to outcomes were
those calling for \emph{repeated} decisions distributed across a stay
(routine basal insulin dosing, routine glucose monitoring, insulin
following hyperglycemia), which were also the ones least well
implemented.

Whether \emph{feeding} such measured signals into a predictive model
improves what the model does has, to our knowledge, not previously been
tested. This paper does so in a controlled fashion, in three of the seven
ladder arms.

\subsection{The comparison design}

Every ladder claim in this paper---architecture, knowledge base, quality
signal---requires that anything else that could shift performance is
identical across the compared arms. The clinical adherence literature
learned this lesson twice over: guideline-outcome associations quietly
fail when eligibility, exposure window, follow-up horizon, or covariate
handling drift across compared groups. The corresponding pitfall in
short-horizon prediction is that each modelling arm can silently see a
different label set, a different patient set, or a different observation
window if its own preprocessing decides them independently. Section~\ref{sec:methods}
describes the shared-cohort machinery that keeps the seven arms on
identical labels and identical patients, and Section~\ref{sec:methods-guards}
describes the leakage and knowledge-containment guards that make each
contrast measure what it is meant to.
